%! From Combinations to Solutions --- Setting Up Ax = b — Unit 2, Lesson 2.0 — Homework
\documentclass[10pt]{article}
\usepackage{linalg-article}
\usepackage{linalg-boxes}
\newcommand{\TallMath}[1]{$\displaystyle #1\rule[-1.4em]{0pt}{3.2em}$}
\newcommand{\xx}{\mathbf{x}}
\newcommand{\bb}{\mathbf{b}}

\begin{document}

\pageheader{Unit 2, Lesson 2.0}{Homework}
\namedateperiod

\vspace{0.12in}

\begin{notesbox}{Practice}
Show your steps and \textbf{check} every solution by rebuilding $\bb$.

\begin{enumerate}[leftmargin=1.9em, itemsep=12pt]
  \item \textbf{Three views.} Let $A=\begin{bmatrix} 2 & 1 \\ 1 & 4 \end{bmatrix}$ and
        $\bb=\begin{bmatrix} 8 \\ 11 \end{bmatrix}$.
    \begin{enumerate}[label=(\alph*), leftmargin=1.8em, itemsep=4pt]
      \item Write $A\xx=\bb$ as a \emph{combination of the columns} of $A$.
      \item Write $A\xx=\bb$ as a \emph{system of two equations}.
    \end{enumerate}
    \vspace{1.3cm}
  \item \textbf{Solve by elimination.} Find $x_1$ and $x_2$ for each system.
    \begin{enumerate}[label=(\alph*), leftmargin=1.8em, itemsep=4pt]
      \item $\begin{aligned}[t] x_1 + x_2 &= 6 \\ 2x_1 + x_2 &= 8 \end{aligned}$ \hfill
      \item $\begin{aligned}[t] 3x_1 + 2x_2 &= 16 \\ x_1 + 2x_2 &= 8 \end{aligned}$ \hfill
      \item $\begin{aligned}[t] x_1 - x_2 &= 1 \\ 2x_1 + x_2 &= 8 \end{aligned}$
    \end{enumerate}
    \vspace{1.6cm}
  \item \textbf{Fertilizer blend.} Each bag of \textbf{Base A} has $1$ kg nitrogen and $2$ kg
        phosphorus; each bag of \textbf{Base B} has $3$ kg nitrogen and $1$ kg phosphorus. A
        field needs $13$ kg nitrogen and $11$ kg phosphorus.
    \begin{enumerate}[label=(\alph*), leftmargin=1.8em, itemsep=6pt]
      \item Set up and solve $A\xx=\bb$ for the number of bags of each base. \vspace{1.2cm}
      \item \emph{Interpret:} how many bags of each, and does the answer make sense? \writeline
    \end{enumerate}
  \item \textbf{How many solutions?} Classify each system as \emph{one}, \emph{none}, or
        \emph{infinitely many}, and say why (parallel lines, same line, or crossing).
    \begin{enumerate}[label=(\alph*), leftmargin=1.8em, itemsep=4pt]
      \item $\begin{aligned}[t] x_1 + 2x_2 &= 5 \\ 2x_1 + 4x_2 &= 10 \end{aligned}$ \hfill
      \item $\begin{aligned}[t] x_1 + 2x_2 &= 5 \\ 2x_1 + 4x_2 &= 7 \end{aligned}$ \hfill
      \item $\begin{aligned}[t] x_1 + 2x_2 &= 5 \\ x_1 + 3x_2 &= 7 \end{aligned}$
    \end{enumerate}
    \vspace{1.2cm}
\end{enumerate}
\end{notesbox}

\begin{extensionbox}
\textbf{One variable at a time.} Solve this $3\times 3$ system by eliminating $x$ first:
\[
  \begin{aligned}
    x + y + z &= 6 \\
    x + 2y + z &= 8 \\
    2x + y + z &= 7
  \end{aligned}
\]
Subtract the first equation from the second, and subtract twice the first from the third. You
are left with a small system in $y$ and $z$ --- solve it, then back-substitute for $x$.
\writelines{3}
\end{extensionbox}

\begin{spiralbox}
\textbf{Coming next --- Lesson 2.1: The Idea of Elimination.} Today you subtracted one equation
from another to make a variable disappear. Next we turn that into a reliable, step-by-step
\emph{procedure} --- choosing \textbf{pivots} and clearing out the entries below them --- so it
works for any size system, not just the lucky ones.
\end{spiralbox}

\end{document}
